\documentclass[12pt]{article}

\usepackage{fouriernc}
\usepackage{amsmath}

\usepackage{tikz}

\newcommand{\contradiction}{%
\begin{tikzpicture}[rotate=45,x=0.5ex,y=0.5ex]
\draw[line width=.2ex] (0,2) -- (3,2) (0,1) -- (3,1) (1,3) -- (1,0) (2,3) -- (2,0);
\end{tikzpicture}
}

\begin{document}

\section*{Problem-183}
\subsection*{Problem Statement}
The first term of an arithmetic progression is $a$, the 4th term is $b$. Find the second and the third terms.
\subsection*{Solution}
The $n$-th term of the arithmetic progression has the form $a+(n-1)d$ where $a$ is the first term ($n=1$) and $d$ is the difference. From the information given in the statement, we have:
\[
a + (4-1)d = b
\]
Thus $d = \frac{b-a}{3}$. So, the second term is $a+(2-1)d = a+\frac{b-a}{3} = \frac{2a+b}{3}$. The third term is $d$ more than the second term, that is, $\frac{a+2b}{3}$.

\section*{Problem-198}
\subsection*{Problem Statement}
The first term of a geometric progression is $a$ and the third term is $b$. Find the second term.
\subsection*{Solution}
Say the second term is $x$. Then from the definition of a geometric progression, $\frac{x}{a} = \frac{b}{x}$. Therefore, $x^2 = ab$. We shall consider three separate possibilities.
\begin{enumerate}
\item When $ab < 0$, there is no real-valued $x$. For example, a geometric progression will all real-valued terms, even if the ratio were negative like -1, the terms of the geometric progression would have alternate signs, so that the first and the third terms would still have the same sign, making $ab < 0$ impossible.
\item When $ab = 0$, we have $x = 0$.
\item When $ab > 0$, there are two possibilities for $x$: $\sqrt{ab}$ and $-\sqrt{ab}$. 
\end{enumerate}

\section*{Problem-204}
\subsection*{Problem Statement}
Is it possible that numbers $\frac{1}{2}$, $\frac{1}{3}$, and $\frac{1}{5}$ are (not necessarily adjacent) terms of the same arithmetic progression?

\subsection*{Solution}
If the numbers come from the same arithmetic progression, then we can consider them as sorted. Since, in the question, the numbers are already in descending order, we shall go with that order. So, if the three numbers indeed come from the same arithmetic progression, the difference will turn out to be negative. We may consider $\frac{1}{2}$ as the first term of the sub-progression. We thus have:
\begin{equation*}
\begin{aligned}
\frac{1}{2} + (m-1)\ d &= \frac{1}{3}\\
\frac{1}{2} + (n-1)\ d &= \frac{1}{5}
\end{aligned}
\end{equation*}
Where $m \neq 1$, $m < n$ and $d\ (\neq 0)$ is the difference. We can rearrange and get the below two equations:
\begin{equation}
		(m-1)\ d = -\frac{1}{6}	
\end{equation}
\begin{equation}
		(n-1)\ d = -\frac{3}{10}	
\end{equation}
Dividing $(1)$ by $(2)$, we have:
\begin{equation}
	\frac{m-1}{n-1} = \frac{5}{9}
\end{equation}
Choosing $(m, n) = (6, 10)$ agrees with $(3)$. Using $m=6$ in $(1)$ gives $d = -\frac{1}{30}$. So, the numbers could be a part of the same arithmetic progression because, $\frac{1}{2} + 5\ \left( -\frac{1}{30} \right) = \frac{1}{3}$ and $\frac{1}{2} + 9 \ \left( -\frac{1}{30} \right) = \frac{1}{5}$. Hence, starting from $\frac{1}{2}$, the 5-th term is $\frac{1}{3}$ and the 9-th term is $\frac{1}{5}$. Actually, there are other valid choices like $(m, n, d) = \left( 11, 19, -\frac{1}{60} \right)$.

\section*{Problem-205}
\subsection*{Problem Statement}
Is it possible that the numbers 2, 3, and 5 are (not necessarily adjacent) terms of a geometric progression?

\subsection*{Solution}
If the numbers belong to the same geometric progression, they must be sorted. Since the numbers are in increasing order in the question, we shall go with that order. So, for the sub-progression that starts at 2, we could consider 2 as the first term. We thus have the following:
\setcounter{equation}{0}
\begin{equation*}
	\begin{aligned}
		2 \cdot q^m &= 3\\
		q^{mn} &= \frac{3^n}{2^n}
	\end{aligned}
\end{equation*}
And,
\begin{equation*}
	\begin{aligned}
		2 \cdot q^n &= 5\\
		     q^{mn} &= \frac{5^m}{2^m}
	\end{aligned}
\end{equation*}
Here $0 < m < n$ and $q$ is the ratio. We can then write:
\begin{equation*}
	\begin{aligned}
		\frac{3^n}{2^n} &= \frac{5^m}{2^m}\\
		3^n &= 2^{n-m} \cdot 5^m
	\end{aligned}
\end{equation*}
On the left side we have an odd number and on the right side we have an even number. We have reached a contradiction. \contradiction\\

Therefore, the numbers cannot be part of the same geometric progression. Note, we would reach a contradiction even if we considered the numbers in reverse order.

\section*{Problem-210}
\subsection*{Problem Statement}
Is it possible that the second term of a geometric progression is less than its first term and also less than its third term?

\subsection*{Solution}
Let's assume that the first three terms of the geometric progression are $a$, $aq$, and $aq^2$. When $a = 0$ or $q = 0$, the answer is negative. So, we shall now consider $a \neq 0$ and $q \neq 0$ case. We shall consider two separate possibilities.
\begin{enumerate}
\item When $a < 0$. If $aq < a$, then $q > 1$. If it is also true that $aq < aq^2$, then $q > q^2$. Both $q > 1$ and $q > q^2$ cannot be true at the same time.
\item When $a > 0$. If $aq < a$, we have $q < 1$. If it is also true that $aq < aq^2$, then $q < q^2$. Both $q < 1$ and $q < q^2$ are possible for say $q = -1$.
\end{enumerate}

So, $1,\ -1,\ 1, \ldots$ is a possible geometric progression.

\section*{Problem-212}
\subsection*{Problem Statement}
Is it possible that an infinite arithmetic progression contains exactly two integer terms?

\subsection*{Solution}
It is not possible. We shall use \textbf{proof by contradiction}. Therefore, we shall assume that exactly two terms of an infinite arithmetic progression are integers and that would lead us to a contradiction, completing the proof.

Say $a+md = x$ and $a+nd = y$ are the only two integral terms where $a$ is the first term of the progression, $0 < m < n$, and $d$ is the difference. We notice that $(n-m)d = y-x$ is an integer. Now consider the term $y + (n-m)d = z$. $z$ is another term distinct from $x$ and $y$. By our assumption, $z$ must not be an integer. However, $y$ and $(n-m)d$ are both integers, so should be their sum, i.e., $z$. We have our contradiction. \contradiction

\section*{Problem-214}
\subsection*{Problem Statement}
In a geometric progression each term is equal to the sum of two preceeding terms. What can be said about the common ratio of this progression?

\subsection*{Solution}
Say $a \neq 0$ is the first term and $q \neq 0$ is the common ratio. Three consecutive terms in the progression would be $aq^n,\ aq^{n+1},$ and $aq^{n+2}$. We thus have:
\begin{equation*}
	\begin{aligned}
		aq^{n+2} &= aq^{n+1} + aq^{n}\\
		q^2 - q - 1 &= 0\\
		q^2 - 2 \cdot q \cdot \frac{1}{2} + \frac{1}{4} - \frac{5}{4} &= 0\\
		\left(q - \frac{1}{2}\right)^2 - \left( \frac{\sqrt{5}}{2} \right)^2 &= 0\\
		\left( q - \frac{1+\sqrt{5}}{2} \right)\ \left( q - \frac{1-\sqrt{5}}{2} \right) &= 0
	\end{aligned}
\end{equation*}
There are two possible values for the common ratio $q$: $\frac{1+\sqrt{5}}{2}$ and $\frac{1-\sqrt{5}}{2}$.

\section*{Problem-215}
\subsection*{Problem Statement}
The \emph{Fibonacci sequence}
\[
	1,\ 1,\ 2,\ 3,\ 5,\ 8,\ 13,\ 21,\ \ldots
\]
is defined as follows: The first two terms are equal to 1, and each subsequent term is equal to the sum of the two preceeding terms. Find numbers $A$ and $B$ such that (for all $n$), the $n$th term of the Fibonacci sequence is equal to
\[
	A \left( \frac{1+\sqrt{5}}{2} \right)^n + B \left( \frac{1-\sqrt{5}}{2} \right)^n
\]

\subsection*{Solution}
For $n=1,2$ the terms of the Fibonacci sequence are both 1. We shall use this observation to derive two equations and we should be able to find $A$ and $B$ as a solution-pair of the two equations. Then we shall use induction to show that those values of $A$ and $B$ once plugged into the formula for the $n$th term, work for all $n$.

\subsubsection*{Find values of $A$ and $B$}
From the first two terms ($n=1,2$) of the Fibonacci sequence, we have the below pair of equations.
\setcounter{equation}{0}
\begin{equation}
	A \left( \frac{ 1+\sqrt{5} }{2} \right) + B \left( \frac{ 1-\sqrt{5} }{2} \right) = 1
\end{equation}
\begin{equation}
	A \left( \frac{ 1+\sqrt{5} }{2} \right)^2 + B \left( \frac{ 1-\sqrt{5} }{2} \right)^2 = 1
\end{equation}
By massaging $(1)$ we get $(3)$ and by massaging $(2)$ we get $(4)$:
\begin{equation}
	A(1+\sqrt{5})^2 -4B = 2(1 + \sqrt{5})
\end{equation}
\begin{equation}
	A(1+\sqrt{5})^2 + B(1 - \sqrt{5})^2 = 4
\end{equation}
From $(4)-(3)$ we get:
\begin{equation*}
	\begin{aligned}
		B(1 - \sqrt{5})^2 + 4B &= 4-2(1 + \sqrt{5})\\
		B &= \frac{ 2(1-\sqrt{5}) }{ -2\sqrt{5}(1-\sqrt{5}) }\\
		B &= -\frac{1}{\sqrt{5}}
	\end{aligned}
\end{equation*}
Plugging in $B = -\frac{1}{\sqrt{5}}$ in $(1)$ gives us:
\begin{equation*}
	\begin{aligned}
		A\left( \frac{1+\sqrt{5}}{2} \right) -\frac{1}{\sqrt{5}}\left( \frac{1-\sqrt{5}}{2} \right) &= 1\\
		A \left( \frac{1+\sqrt{5}}{2} \right) &= \frac{ 1+\sqrt{5} }{ 2\sqrt{5} }\\
		A &= \frac{1}{\sqrt{5}}
	\end{aligned}
\end{equation*}
So, the potential formula for the $n$th term of the Fibonacci sequence is as follows:
\[
	\frac{ 1 }{ \sqrt{5} } \left( \frac{ 1+\sqrt{5} }{2} \right)^n - \frac{1}{ \sqrt{5} } \left( \frac{1-\sqrt{5}}{2} \right)^n
\]

\subsubsection*{Prove that the formula works for all $n$}
We already know that the formula works for $n=1,2$. We shall use these as our base cases. For $n > 2$ we shall use induction. Therefore, we shall show that adding $n$th and $(n+1)$th terms yields the $(n+2)$th term. For convenience, let's denote $\frac{ 1+\sqrt{5} }{2}$ by $\alpha$ and $\frac{ 1-\sqrt{5} }{2}$ by $\beta$. We make the below two observations:
\[
	\alpha^2 = 1+\alpha
\]
\[
	\beta^2 = 1+\beta
\]
Now to complete induction, we add up $n$th and $(n+1)$th terms which should give us the $(n+2)$th term.
\begin{equation*}
	\begin{aligned}
		& \left( \frac{1}{\sqrt{5}}\ \alpha^{n+1} - \frac{1}{\sqrt{5}}\ \beta^{n+1} \right) + \left( \frac{1}{\sqrt{5}}\ \alpha^n - \frac{1}{\sqrt{5}}\ \beta^n \right)\\
		&= \frac{\alpha^n}{\sqrt{5}}\ (1+\alpha) - \frac{\beta^n}{\sqrt{5}}\ (1+\beta)\\
		&= \frac{\alpha^n}{\sqrt{5}}\ \alpha^2 - \frac{\beta^n}{\sqrt{5}}\ \beta^2\\
		&= \frac{1}{\sqrt{5}}\ \alpha^{n+2} - \frac{1}{\sqrt{5}}\ \beta^{n+2}\ \ \ \blacksquare
	\end{aligned}
\end{equation*}

\end{document}